\chapter{Partial parametrization}\label{app-partial}

Section~\ref{sec-proc-lifetimes} showed that an Algol 68 function cannot
be returned from the block whose names it uses: a function that composes
two other functions, the everyday closure of current languages, is an
error. Soon after the Revised Report, Lindsey proposed an extension that
removes most of this limitation without changing the lifetime
rule~\cite{partial}. The subcommittee of IFIP Working Group 2.1 that
looked after Algol 68 released it in 1976, so that implementations that
experimented with such a facility would all provide it in the same form;
it is not part of the Revised Report. Algol 68 Genie implements it;
ga68 and a68j do not (Appendix~\ref{app-impl}): a68j rejects such a call
with a message, and ga68 currently stops with an internal compiler
error. The programs of this
appendix were run with Algol 68 Genie.

\section*{Leaving arguments out}

With partial parametrization, a call may leave arguments out, by writing
nothing in their places; the commas stay. Such a call does not run the function; it
yields a new function that takes the missing arguments and has the given
ones stored inside it:

\example{appG/add}

\noindent \code|add (3, )| is a function of one \code|real|, which adds 3.
It can be called at once, as in \code|add (3, ) (10)|, or kept and called
later. The functional languages call the same idea partial application;
in Python it is \texttt{functools.partial}, in C++ \texttt{std::bind}.

A function can be given its arguments in several stages. Each stage
stores the new arguments and yields a function of the remaining ones:

\example{appG/stages}

\section*{Closures without the lifetime problem}

The stored arguments are part of the new function: they are copied into
it, and the function does not refer to the block in which the call was
made. This is what makes the idea work. Instead of using the names of an
enclosing block, the inner function receives everything as parameters,
and the outer function supplies them by partial parametrization. The
function \code|compose| of Section~\ref{sec-proc-lifetimes}, which Algol
68 Genie stops at run time there, becomes:

\example{appG/compose}

\noindent The inner function uses only its own parameters
\code|f|, \code|g| and \code|x|. \code|compose| gives it the first two and
returns the result, which no longer depends on \code|compose|'s
parameters. Lindsey's proposal shows the same pattern with an operator
that raises the value of a function to a power, so that the cosine
``squared'' yields $\cos^2\theta$. The operator \code|^| below instead
applies a function \code|n| times:

\example{appG/power}

\noindent Twenty applications of \code|cos| to 1 come close to the
number $x$ with $\cos x = x$, 0.739085\ldots; more applications get
closer.

\section*{When the arguments are evaluated}

The arguments of a call that leaves arguments out are evaluated when the call is made, not
when the resulting function is finally called. A value is stored as it
is at that moment; a reference is stored as a reference, so the function
sees later changes of the variable it refers to:

\example{appG/when}

\noindent \code|add_n| keeps the value 1 that \code|n| had when
\code|add_n| was made; the later assignment to \code|n| does not change
it. \code|bump| keeps the reference \code|counter|, and every call
increases the same cell. A heap cell and a function that keeps a
reference to it are the counter of Section~\ref{sec-proc-lifetimes}
without the need to pass the counter along.

\begin{pitfall}
The lifetime rule still applies to what is stored. A function made by
partial parametrization lives only as long as both the original function
and every value stored in it, whichever ends first. If a stored reference refers to a \code|loc| cell of a
block, the function must not outlive that block:

\example{appG/scope}

\noindent The cell \code|count| ends with the block of
\code|make_bump|, and the function returned holds a reference to it. By
the rules of the proposal, returning the function from the block is the
error. Algol 68 Genie does not check it there; it stops the program only
when the function is used, here in \code|incr|, whose result is the cell
that no longer exists. The counter must be a \code|heap| cell, as in the
previous example.
\end{pitfall}

\section*{Calls with all arguments}

In Lindsey's proposal, every call is a partial parametrization: a call
that supplies all arguments yields a function without parameters, which
is then called at once if the context requires a value, the usual case.
If the context requires a function without parameters, the call is not
executed but kept: \code|proc int tick = incr (counter, 1)| would declare
a function that increases \code|counter| by 1 each time it is called.

\begin{implnote}
Algol 68 Genie does not implement this part of the proposal. It runs a
call as soon as all arguments are given, and rejects the declaration of
\code|tick|:
\outputof{appG/complete}
\noindent Leaving out at least one argument, as \code|bump| does above,
works in Algol 68 Genie.
\end{implnote}
